\documentclass[11pt]{article}
\usepackage[margin=1in]{geometry}
\usepackage{amsmath}
\usepackage{amssymb}
\usepackage{enumitem}
\usepackage{hyperref}

\title{Summary of Branch Changes and Timestep Constraints}
\author{}
\date{}

\begin{document}
\maketitle

\section{Branch Summary Relative to \texttt{develop}}

This branch adds an optional CFL-like contact stiffness timestep limit on top of the existing Common Plane timestep vote. The main code changes are:

\begin{itemize}[leftmargin=*]
\item A new API hook for enabling the stability-limit path.
\item A small refactor that exposes the uncoupled penalty stiffness calculation used by Common Plane penalty contact.
\item A third timestep-vote bucket in \texttt{CouplingScheme::computeCommonPlaneTimeStep()}, used for the CFL-like stiffness limit.
\item A unit test that verifies the implemented CFL-like limit.
\end{itemize}

The force law itself is still the Common Plane penalty law,
\begin{equation}
  p = g \, k_c,
\end{equation}
where $g$ is the signed normal gap and $k_c$ is the equivalent interface penalty stiffness per area. In this branch,
\begin{equation}
  k_c = \frac{k_1 k_2}{k_1 + k_2},
\end{equation}
with
\begin{equation}
  k_i =
  \begin{cases}
    s_i k_{\mathrm{const},i}, & \text{constant penalty}, \\
    s_i K_i / t_i, & \text{element-based penalty},
  \end{cases}
\end{equation}
where $s_i$ is the mesh penalty scale, $K_i$ is the registered bulk modulus, and $t_i$ is the registered element thickness.

\section{Checks Implemented in the Timestep Vote}

The branch now returns the minimum over three groups of checks:
\begin{equation}
  \Delta t_{\mathrm{vote}}
  = \min\!\left(
      \Delta t_{\mathrm{gap}},
      \Delta t_{\mathrm{proj}},
      \Delta t_{\mathrm{stab}}
    \right).
\end{equation}

\subsection*{Existing kinematic checks}

\textbf{1. Current-gap limit.}
For active contact, the code compares the current interpenetration against an allowable fraction of each side's thickness,
\begin{equation}
  \delta_i^{\max} = \alpha_p t_i,
\end{equation}
with $\alpha_p = \texttt{timestep\_pen\_frac}$. If the current overlap already exceeds this bound and the projected motion continues to drive penetration, the vote is
\begin{equation}
  \Delta t_i = \alpha \frac{\delta_i^{\max}}{v_{i,n_i}},
\end{equation}
using the face-normal velocity projection and the user scale factor $\alpha = \texttt{timestep\_scale}$.

\textbf{2. Velocity-projection limit.}
Even for candidate pairs that are not currently active, the overlap centroid is advanced by one host step using the interpolated face velocities. If that projected penetration exceeds $\delta_i^{\max}$, the same reset-gap estimate is used:
\begin{equation}
  \Delta t_i = \alpha \frac{\delta_i^{\max}}{v_{i,n_i}}.
\end{equation}

These two checks are geometry and kinematics based. They are intended to avoid excessive penetration and missed contact resolution, not to estimate the dynamic stability limit of the penalty spring.

\subsection*{New stability check}

The new limit is applied only when the pair is in active contact and \texttt{enable\_timestep\_stability\_limits} is true.

\textbf{3. CFL-like stiffness limit.}
The branch rescales the host timestep using the added contact stiffness:
\begin{equation}
  f_i = \frac{1}{\sqrt{1 + k_c/k_i}},
  \qquad
  \Delta t_{\mathrm{CFL}} = \alpha \, \Delta t_{\mathrm{host}} \min(f_1,f_2).
\end{equation}

\[
  \Delta t_{\mathrm{stab}} = \Delta t_{\mathrm{CFL}}.
\]

\section{Does the Remaining Constraint Make Sense?}

\subsection*{What fits well}

\begin{itemize}[leftmargin=*]
\item The CFL-like limit is built from the same equivalent interface stiffness $k_c$ that the penalty force law already uses. That is the right stiffness quantity to reference for Common Plane penalty contact.
\item Reusing the uncoupled side stiffnesses $k_1$ and $k_2$ is consistent with the branch's spring-in-series model for the two contacting bodies.
\item The CFL-like limit is a sensible extension of the existing vote. Tribol does not know the host's lumped masses or modal content, so rescaling the host timestep by the stiffness increase is a pragmatic approximation.
\item Keeping the original gap-based checks is still useful. They control contact resolution quality and candidate detection in ways the stiffness-based limit does not.
\end{itemize}

\subsection*{Remaining caveat}

\begin{itemize}[leftmargin=*]
\item The CFL-like limit is still a host-timestep rescaling, not a true mass-based critical timestep. It is best viewed as a conservative approximation tied to the added penalty stiffness.
\end{itemize}

\subsection*{Bottom line}

For the implemented \emph{Common Plane + penalty} algorithm, the branch remains directionally sound:
\begin{itemize}[leftmargin=*]
\item The old two checks remain appropriate for limiting penetration.
\item The CFL-like stability limit is the most defensible dynamic addition because it uses the same interface stiffness model as the force calculation.
\end{itemize}

\section{Test Coverage}

The branch adds a direct formula check for the CFL-like limit. That test confirms that the code returns the implemented formula. It does \emph{not} independently validate the physical derivation of the host-timestep rescaling.

\section{Practical Recommendation}

If the goal is a conservative host-code timestep vote for explicit Common Plane penalty contact, the remaining CFL-like limit is reasonable to keep. If the goal is a physically grounded critical timestep estimate, the next step should be to derive a mass-aware bound from the actual semi-discrete contact model.

\end{document}
